% ch10_c 第10章 §10.4尾–§10.6.2 (书页 466–478 / p481–p493)
% 块界说明：
%   起点——书页 466（p481）顶部 "fermion PSGs that we obtained above. ..." 是
%   书页 465（p480，ch10_b 块）末句 "The PSGs obtained in Section 10.2 ... agree
%   exactly with the" 的延续。ch10_b 已把前半句译文（"……得到的 PSG 与"）留在
%   其正文最后，未设 TAIL 区。本块正文第一个非注释行以 %JOIN% 起头续译
%   "我们在上面得到的费米子 PSG 完全一致。……"，不重复前半句。
%   终点——书页 478（p493）末段以 "...around the i site."（"……绕 i 格点的
%   两个最近邻转子。"）整句自然收束；书页 479（p494）顶部 "In the following,
%   we are going to show that..." 为新段，属下一块（ch10_d），故本块自然结束，
%   不设 %--C-TAIL-- 区。
%   另：蓝本把问题 10.5.2 的副题切成伪标题 "## Solving eqn (10.5.4) ..."，
%   按 PNG（p483）原书为题内粗体行，已还原为 \textbf 行。
%JOIN%我们在上面得到的费米子 PSG 完全一致。这个例子表明，9.4 节中引入的、用于描述量子序的那些 PSG，正是凝聚弦端点的跳跃哈密顿量的对称性群。量子序的 PSG 描述与弦网凝聚描述是密切相关的。

这里我们想指出，9.4 节中引入的 PSG 都是费米子 PSG。它们只是可用于刻画量子序的众多不同种类的 PSG 中的一种。一般地说，一个量子有序态可以包含几种不同类型弦的凝聚，每一类凝聚弦的端点都将拥有各自的 PSG。找出所有类型的凝聚弦及其 PSG，将使我们能够对量子序作出更完备的刻画。

\section{立方格子上的演生费米子与弦网凝聚}

\begin{keypoints}
\item 三维模型中的演生费米子，无法用把通量附着到玻色子上的方式来理解。要理解三维中的演生费米子，确实需要弦网理论。
\end{keypoints}

在 $2+1$ 维中，把 $\pi$ 通量束缚到一个单位电荷上（或者把一个 $Z_2$ 涡旋束缚到一个 $Z_2$ 电荷上），便可以得到一个费米子。在我们的自旋 $1/2$ 模型 $H_U+H_g$ 中，我们正是这样得到费米子的。由于 $Z_2$ 涡旋与 $Z_2$ 电荷都以开弦端点的形式出现，费米子也以弦端点的形式出现。我们看到，对于 $(2+1)$ 维玻色模型中的演生费米子，我们有两套理论：演生费米子既可以视为电荷与通量的束缚态，也可以视为弦的端点。然而，在 $3+1$ 维中，我们不能通过附着 $\pi$ 通量把玻色子变成费米子。于是人们或许会问：费米子能否从 $(3+1)$ 维的局域玻色模型中演生出来？本节我们将研究立方格子上一个精确可解的自旋 $3/2$ 模型。我们将研究这样一个模型中的弦网凝聚态，并证明费米子仍然可以作为弦的端点出现。事实上，关于演生费米子的弦网理论在任意维数下都成立。

\subsection{立方格子上的精确可解自旋 $3/2$ 模型}

为了在三维立方格子上构造一个精确可解模型，我们先引入算符 $\gamma_{\pmb{i}}^{ab}$，$a,b=x,\bar{x},y,\bar{y},z,\bar{z}$，它们作用在每个格点的局域希尔伯特空间之内。在一个格点上，$\gamma^{ab}$ 满足以下代数（格点指标 $\pmb{i}$ 已略去）：
\begin{align*}
\gamma^{ab}&=-\gamma^{ba}=(\gamma^{ab})^{\dagger}, & (\gamma^{ab})^{2}&=1, \tag{10.5.1}\\
[\gamma^{ab},\gamma^{cd}]&=0, & \gamma^{ab}\gamma^{bc}&=\mathrm{i}\gamma^{ac}, \qquad \text{若 } a,b,c,d \text{ 互不相同}
\end{align*}

为了证明这个代数是自洽的，令 $\lambda^{a}$ 为一个由 $a=x,\bar{x},y,\bar{y},z,\bar{z}$ 标记的马约拉纳费米子算符。从马约拉纳费米子的代数 $\{\lambda^{a},\lambda^{b}\}=2\delta_{ab}$ 出发可以证明，
\begin{equation*}
\gamma^{ab}=\frac{\mathrm{i}}{2}\left(\lambda_{\pmb{i}}^{a}\lambda_{\pmb{i}}^{b}-\lambda_{\pmb{i}}^{b}\lambda_{\pmb{i}}^{a}\right)
\end{equation*}
满足上述代数。

上述代数有一个四维表示。我们首先注意到，$\gamma^{az}\equiv\gamma^{a}$（$a=x,\bar{x},y,\bar{y}$）满足狄拉克代数
\begin{equation*}
\{\gamma^{a},\gamma^{b}\}=2\delta_{ab}, \qquad a,b=x,\bar{x},y,\bar{y}. \tag{10.5.2}
\end{equation*}
因此，我们可以用四个 $4\times4$ 狄拉克矩阵把它们表示为
\begin{align*}
\gamma^{xz}=\gamma^{x}=\sigma^{x}\otimes\sigma^{x}, & & \gamma^{\bar{x}z}=\gamma^{\bar{x}}=\sigma^{y}\otimes\sigma^{x}, &\\
\gamma^{yz}=\gamma^{y}=\sigma^{z}\otimes\sigma^{x}, & & \gamma^{\bar{y}z}=\gamma^{\bar{y}}=\sigma^{0}\otimes\sigma^{y}. & \tag{10.5.3}
\end{align*}
我们还可以定义
\begin{equation*}
\gamma^{5}=\gamma^{x}\gamma^{\bar{x}}\gamma^{y}\gamma^{\bar{y}}=-\sigma^{0}\otimes\sigma^{z}
\end{equation*}
它与 $\gamma^{a}$ 反对易，即 $\{\gamma^{a},\gamma^{5}\}=0$。由于 $\gamma^{z\bar{z}}$ 与 $\gamma^{az}\equiv\gamma^{a}$（$a=x,\bar{x},y,\bar{y}$）反对易，我们可以把 $\gamma^{5}$ 认同为 $\gamma^{z\bar{z}}$。利用代数 (10.5.1)，可以把其余的 $\gamma^{ab}$ 表示为
\begin{equation*}
\gamma^{a\bar{z}}=\mathrm{i}\gamma^{a}\gamma^{5}, \qquad a=x,\bar{x},y,\bar{y}
\end{equation*}
这样，我们就把全部 $\gamma^{ab}$（$a,b=x,\bar{x},y,\bar{y},z,\bar{z}$）用狄拉克矩阵 $\gamma^{x,\bar{x},y,\bar{y}}$ 表示了出来。可以证明，这样定义的 $\gamma^{ab}$ 满足代数 (10.5.1)。

由于 $\gamma^{ab}$ 是四维的，我们可以说它们作用在自旋 $3/2$ 态上。借助 $\gamma_{\pmb{i}}^{ab}$，我们可以在立方格子上写出一个精确可解的自旋 $3/2$ 模型如下：
\begin{align*}
H_{3D}&\\
&=g\sum_{\pmb{i}}\Big(\gamma_{\pmb{i}}^{yx}\gamma_{\pmb{i}+\pmb{x}}^{\bar{x}y}\gamma_{\pmb{i}+\pmb{x}+\pmb{y}}^{\bar{y}\bar{x}}\gamma_{\pmb{i}+\pmb{y}}^{x\bar{y}}+\gamma_{\pmb{i}}^{zy}\gamma_{\pmb{i}+\pmb{y}}^{\bar{y}z}\gamma_{\pmb{i}+\pmb{y}+\pmb{z}}^{\bar{z}\bar{y}}\gamma_{\pmb{i}+\pmb{z}}^{y\bar{z}}+\gamma_{\pmb{i}}^{xz}\gamma_{\pmb{i}+\pmb{z}}^{\bar{z}x}\gamma_{\pmb{i}+\pmb{z}+\pmb{x}}^{\bar{x}\bar{z}}\gamma_{\pmb{i}+\pmb{x}}^{z\bar{y}}\Big)\\
&=-g\sum_{\pmb{p}}\hat F_{\pmb{p}} \tag{10.5.4}
\end{align*}
其中 $\pmb{p}$ 标记立方格子中的方格，而视方格 $\pmb{p}$ 的取向，$\hat F_{\pmb{p}}$ 等于 $-\gamma_{\pmb{i}}^{yx}\allowbreak\gamma_{\pmb{i}+\pmb{x}}^{\bar{x}y}\allowbreak\gamma_{\pmb{i}+\pmb{x}+\pmb{y}}^{\bar{y}\bar{x}}\allowbreak\gamma_{\pmb{i}+\pmb{y}}^{x\bar{y}}$、$-\gamma_{\pmb{i}}^{zy}\allowbreak\gamma_{\pmb{i}+\pmb{y}}^{\bar{y}z}\allowbreak\gamma_{\pmb{i}+\pmb{y}+\pmb{z}}^{\bar{z}\bar{y}}\allowbreak\gamma_{\pmb{i}+\pmb{z}}^{y\bar{z}}$ 或 $-\gamma_{\pmb{i}}^{xz}\allowbreak\gamma_{\pmb{i}+\pmb{z}}^{\bar{z}x}\allowbreak\gamma_{\pmb{i}+\pmb{z}+\pmb{x}}^{\bar{x}\bar{z}}\allowbreak\gamma_{\pmb{i}+\pmb{x}}^{z\bar{y}}$。

注意到
\begin{equation*}
[\hat F_{\pmb{p}},\hat F_{\pmb{p}}']=0, \qquad (\hat F_{\pmb{p}})^{2}=1,
\end{equation*}
就容易求出 $H_{3D}$ 的本征值。设 $|\{F_{\pmb{p}}\},\alpha\rangle$ 是 $\hat F_{\pmb{p}}$ 的共同本征态，即 $\hat F_{\pmb{p}}|\{F_{\pmb{p}}\},\alpha\rangle=F_{\pmb{p}}|\{F_{\pmb{p}}\},\alpha\rangle$，其中 $\alpha$ 标记可能的简并。那么 $|\{F_{\pmb{p}}\},\alpha\rangle$ 也是能量本征态，能量为 $-g\sum_{\pmb{p}}F_{\pmb{p}}$。由于 $F_{\pmb{p}}=\pm1$，在周期性格子上基态能量为 $E_0=-3|g|N_{\text{site}}$，其中 $N_{\text{site}}$ 是格点数目。激发态可以通过翻转某些 $F_{\pmb{p}}$ 的符号得到，其能量为 $E_0$ 加上 $2|g|$ 的整数倍。

有一个微妙之处：诸 $F_{\pmb{p}}$ 并不全部独立。若 $S$ 是一个单胞立方体的表面，则有算符恒等式
\begin{equation*}
\prod_{\pmb{p}\in S}\hat F_{\pmb{p}}=1
\end{equation*}
由此可知，对所有立方体都有 $\prod_{\pmb{p}\in S}F_{\pmb{p}}=1$。如果我们把 $F_{\pmb{p}}$ 看作穿过方格 $\pmb{p}$ 的 $Z_2$ 通量，那么上述约束就意味着通量守恒。这意味着我们模型的能谱与立方格子上一个 $Z_2$ 规范理论的能谱相同。基态可以看作无通量的态，即对所有 $\pmb{p}$ 都有 $F_{\pmb{p}}=1$。基态之上的激发由通量圈描述。基本激发对应于最小的通量圈，即在与某条链 $\langle\pmb{i}\pmb{j}\rangle$ 相邻的四个方格 $\pmb{p}$ 上 $F_{\pmb{p}}=-1$。我们可以把这些激发看作生活在立方格子\emph{链}上的准粒子（见图 10.7）。正如我们后面将看到的，这种通量圈激发具有费米统计。

\subsection*{问题 10.5.1}

\begin{enumerate}
\item 证明式 (10.5.3) 中的 $4\times4$ 矩阵 $\gamma^{a}$ 满足狄拉克代数。
\item 证明式 (10.5.4) 求和中的各项彼此对易。（提示：可以利用 $\gamma^{ab}$ 的马约拉纳费米子表示。）这使得我们能够精确求解 $H_{3D}$。
\end{enumerate}

\subsection*{问题 10.5.2}

\textbf{用马约拉纳费米子和投影构造求解式 (10.5.4)}

\begin{enumerate}
\item 设 $\lambda_{\pmb{i}}^{a}$ 是马约拉纳费米子算符，其中 $a=x,\bar{x},y,\bar{y},z,\bar{z}$。对每条最近邻链，定义
\begin{equation*}
\hat U_{\pmb{i}\pmb{j}}=\lambda_{\pmb{i}}^{a}\lambda_{\pmb{j}}^{b} \tag{10.5.5}
\end{equation*}
其中 $ab$ 是与链 $\pmb{i}\to\pmb{j}$ 相联系的指标对（见图 10.7）。证明诸 $\hat U_{\pmb{i}\pmb{j}}$ 构成一组相互对易的算符。
\item 考虑费米子哈密顿量
\begin{equation*}
H=-\sum_{\pmb{p}}g\hat F_{\pmb{p}}, \qquad \hat F_{\pmb{p}}=\hat U_{\pmb{i}_1,\pmb{i}_2}\hat U_{\pmb{i}_2,\pmb{i}_3}\hat U_{\pmb{i}_3,\pmb{i}_4}\hat U_{\pmb{i}_4,\pmb{i}_1} \tag{10.5.6}
\end{equation*}
%FIG{10.7}: {立方格子中的一条开弦，以及与最近邻链相联系的指标对。开弦的一个端点翻转四个阴影方格上的 $F_{\pmb{p}}$ 的符号。}
其中 $\sum_{\pmb{p}}$ 对立方格子的所有方格面求和。这里 $\pmb{i}_1$、$\pmb{i}_2$、$\pmb{i}_3$、$\pmb{i}_4$ 标记方格 $\pmb{p}$ 的四个角。引入复费米子算符
\begin{equation*}
2\psi_{1,\pmb{i}}=\lambda_{\pmb{i}}^{x}+\mathrm{i}\lambda_{\pmb{i}}^{\bar{x}}, \qquad 2\psi_{2,\pmb{i}}=\lambda_{\pmb{i}}^{y}+\mathrm{i}\lambda_{\pmb{i}}^{\bar{y}}, \qquad 2\psi_{3,\pmb{i}}=\lambda_{\pmb{i}}^{z}+\mathrm{i}\lambda_{\pmb{i}}^{\bar{z}},
\end{equation*}
并令 $N_{\pmb{i}}=\sum_{a=1,2,3}\psi_{a,\pmb{i}}^{\dagger}\psi_{a,\pmb{i}}$ 为格点 $\pmb{i}$ 上的费米子数算符。证明：在每个格点费米子数目为偶的子空间内，费米子哈密顿量 (10.5.6) 约化为自旋 $3/2$ 哈密顿量 (10.5.4)。
\item 用 10.2 节中的程序，求式 (10.5.4) 在具有周期性边界条件的偶数 $\times$ 偶数 $\times$ 偶数格子上的基态能量与基态简并度。
\end{enumerate}

\subsection{弦算符与闭合弦凝聚}

为了证明 $H_{3D}$ 的基态具有闭合弦凝聚，我们需要找到一个与 $H_{3D}$ 对易的闭合弦算符。首先，让我们构造更普遍的开弦算符。

为了用物理的 $\gamma^{ab}$ 算符构造弦算符，我们注意到，可以给每条最近邻链 $\pmb{i}\to\pmb{j}$ 联系一对指标 $ab$，如图 10.7 所示。例如，链 $\pmb{i}\to\pmb{i}+\pmb{x}$ 对应 $x\bar{x}$，链 $\pmb{i}\to\pmb{i}-\pmb{x}$ 对应 $\bar{x}x$，链 $\pmb{i}\to\pmb{i}+\pmb{y}$ 对应 $y\bar{y}$，等等。开弦
%FIG{10.8}: {回路 $C$ 可以分成两个回路 $C_1$ 与 $C_2$。}
算符则定义为
\begin{equation*}
W(C)=-\left(-\mathrm{i}\gamma_{\pmb{i}_1}^{a_1b_1}\right)\left(-\mathrm{i}\gamma_{\pmb{i}_2}^{a_2b_2}\right)\dots\left(-\mathrm{i}\gamma_{\pmb{i}_n}^{a_nb_n}\right) \tag{10.5.7}
\end{equation*}
其中 $b_ma_{m+1}$ 是与链 $\pmb{i}_m\to\pmb{i}_{m+1}$ 相联系的指标对。例如，图 10.7 中那条开弦的开弦算符由下式给出：
\begin{equation*}
W(C)=-\left(-\mathrm{i}\gamma_{\pmb{i}}^{\bar{z}y}\right)\left(-\mathrm{i}\gamma_{\pmb{i}+\pmb{y}}^{\bar{y}x}\right)\left(-\mathrm{i}\gamma_{\pmb{i}+\pmb{x}+\pmb{y}}^{\bar{x}x}\right)
\end{equation*}
我们想强调，开弦 $C$ 是通过相邻格点把诸链的中点连接起来而构成的（见图 10.7）。因此，开弦的端点生活在链上。

闭合弦算符仍由式 (10.5.7) 给出，只是此时 $\pmb{i}_1$ 与 $\pmb{i}_n$ 为最近邻，并且 $b_na_1$ 是与链 $\pmb{i}_n\to\pmb{i}_1$ 相联系的指标对。

让我们先讨论导致闭合弦凝聚的闭合弦算符的一些特殊性质。利用 $\gamma^{ab}$ 的马约拉纳表示，可以证明任意两个闭合弦算符彼此对易。由于 $\hat F_{\pmb{p}}$ 本身就是一个闭合弦算符，闭合弦算符与自旋 $3/2$ 哈密顿量对易。因此，基态也是闭合弦算符的本征态，具有闭合弦凝聚（或更一般地，闭合弦网凝聚）。

如果我们把一个回路 $C$ 分成两个回路 $C_1$ 与 $C_2$（如图 10.8 那样），那么三个回路的闭合弦算符满足关系
\begin{equation*}
W(C)=W(C_1)W(C_2) \tag{10.5.8}
\end{equation*}
这一关系使我们能把 $W(C)$ 写成诸 $F_{\pmb{p}}$ 的乘积，从而计算凝聚闭合弦的振幅。对 $\hat F_{\pmb{p}}=1$ 的基态，我们发现
\begin{equation*}
\langle W(C_{\text{close}})\rangle=1
\end{equation*}
对 $\hat F_{\pmb{p}}=-1$ 的基态，我们发现
\begin{equation*}
\langle W(C_{\text{close}})\rangle=(-)^{N_{\pmb{p}}}
\end{equation*}
其中 $N_{\pmb{p}}$ 是曲面 $S_{C_{\text{close}}}$ 上方格的数目，而 $S_{C_{\text{close}}}$ 是由立方体的面组成、以闭合弦 $C_{\text{close}}$ 为边界的曲面。

\subsection*{问题 10.5.3}

\begin{enumerate}
\item 利用 $\gamma^{ab}=\mathrm{i}\lambda^{a}\lambda^{b}$ 的马约拉纳费米子表示，证明闭合弦算符 (10.5.7) 可以表示为
\begin{equation*}
W(C_{\text{close}})=\hat U_{\pmb{i}_1\pmb{i}_2}\hat U_{\pmb{i}_2\pmb{i}_3}\dots\hat U_{\pmb{i}_{n-1}\pmb{i}_n}\hat U_{\pmb{i}_n\pmb{i}_1}
\end{equation*}
其中 $\hat U_{\pmb{i}\pmb{j}}$ 的定义见式 (10.5.5)。
\item 证明对两个闭合回路定义的闭合弦算符 (10.5.7) 彼此对易。
\item 验证关系式 (10.5.8)。
\end{enumerate}

\subsection{作为开弦端点的演生费米子}

在闭合弦凝聚态中，开弦的端点成为一种新的准粒子。把开弦算符作用到一个基态上，通过计算 $F_{\pmb{p}}$ 与开弦算符之间的对易子可以发现，开弦算符只在其两个端点附近翻转 $F_{\pmb{p}}$ 的符号。因此，弦本身不耗费能量，只有弦端点耗费能量。

注意到弦端点生活在链上，而每条链连接着四个方格，我们发现开弦算符在这四个方格上翻转 $F_{\pmb{p}}$ 的符号（见图 10.7）。因此，开弦的每个端点对应一个小 $Z_2$ 通量圈，耗费能量 $4|2g|=8|g|$。这些小 $Z_2$ 通量圈对应于基态之上的准粒子激发。

为了求出 $Z_2$ 通量圈的统计性质，我们需要用到 Levin 与 Wen (2003) 引入的以下统计跳跃代数（见 4.1.4 节）：
\begin{align*}
&t_{jl}t_{kj}t_{ji}=\mathrm{e}^{\mathrm{i}\theta}t_{ji}t_{kj}t_{jl},\\
&[t_{ij},t_{kl}]=0, \qquad \text{若 } i,j,k,l \text{ 互不相同}, \tag{10.5.9}
\end{align*}
其中 $t_{ji}$ 描述粒子从链 $i$ 到链 $j$ 的跳跃（注意，这里弦的端点生活在链上）。已经证明，若粒子的跳跃算符满足 $\theta=\pi$ 的代数 (10.5.9)，则这些粒子是费米子。如果我们取 $(i,j,k,l)$
%FIG{10.9}: {链 1 上的一个费米子可以跳跃到十条不同的链 2--11 上。跳跃 $1\to4$ 由 $\gamma_{\pmb{i}}^{zx}$ 产生，跳跃 $1\to3$ 由 $\gamma_{\pmb{i}}^{\bar{x}x}$ 产生，跳跃 $1\to7$ 由 $\gamma_{\pmb{j}}^{y\bar{x}}$ 产生。}
为图 10.9 中的链 $(2,1,3,4)$，那么跳跃算符 $(t_{jl},t_{kj},t_{ji})$ 由 $(\gamma_{\pmb{i}}^{xz},\gamma_{\pmb{i}}^{\bar{x}x},\gamma_{\pmb{i}}^{x\bar{z}})$ 给出。式 (10.5.9) 成为
\begin{equation*}
\gamma_{\pmb{i}}^{xz}\gamma_{\pmb{i}}^{\bar{x}x}\gamma_{\pmb{i}}^{x\bar{z}}=\mathrm{e}^{\mathrm{i}\theta}\gamma_{\pmb{i}}^{x\bar{z}}\gamma_{\pmb{i}}^{\bar{x}x}\gamma_{\pmb{i}}^{xz}
\end{equation*}
利用 $\gamma^{ab}$ 的马约拉纳费米子表示，我们发现上式在 $\theta=\pi$ 时成立。如果我们取 $(i,j,k,l)$ 为链 $(2,1,3,10)$，则式 (10.5.9) 成为
\begin{equation*}
\gamma_{\pmb{j}}^{\bar{x}x}\gamma_{\pmb{i}}^{\bar{x}x}\gamma_{\pmb{i}}^{x\bar{z}}=\mathrm{e}^{\mathrm{i}\theta}\gamma_{\pmb{i}}^{x\bar{z}}\gamma_{\pmb{i}}^{\bar{x}x}\gamma_{\pmb{j}}^{\bar{x}x}
\end{equation*}
上式同样在 $\theta=\pi$ 时成立。$(i,j,k,l)$ 的其他取法都属于上述两种类型。弦端点的跳跃满足费米子跳跃代数。因此，弦的端点（即小 $Z_2$ 通量圈）是费米子。我们看到，只要玻色模型具有某些类型闭合弦的凝聚，费米子就能从局域玻色模型中演生出来。

\section{量子转子模型与 $U(1)$ 格点规范理论}

\begin{keypoints}
\item 在更高维数，规范理论相当不平凡。它描述一个纠缠态的涨落，而这些涨落无法用局域自由度来表示。
\item $U(1)$ 规范涨落是空间中闭合弦网的涨落。
\item 如果基态包含任意尺寸闭合弦网的凝聚，就可以演生出无能隙的 $U(1)$ 规范玻色子。
\end{keypoints}

我们已经研究过具有演生 $Z_2$ 规范场的精确可解局域玻色模型。本节我们将从量子转子模型（quantum rotor model）出发，研究一个连续 $U(1)$ 规范理论的演生（Foerster et al., 1980; Senthil and Motrunich, 2002; Wen, 2003a）。
%FIG{10.10}: {(a) 四转子系统。(b) 一个简单的格点规范理论，由格点规范场 $a_{i,i+1}$ 与 $a_0(i)$，$i=1,2,3,4$ 描述。}
我们还将看到，弦网凝聚及相关的量子序如何与演生的 $U(1)$ 规范玻色子相联系。

本节也可以看作是对量子 $U(1)$ 规范理论的一种物理描述。在我看来，用规范势来描述规范理论的标准方式是不自然的，因为规范势本身是非物理的。这里我们将证明，可以用物理自由度——闭合弦网——来表述 $U(1)$ 规范理论（Banks et al., 1977; Savit, 1980）。

\subsection{四转子系统}

\begin{keypoints}
\item 转子角的某些组合的强量子涨落导致一个低能有效规范理论。
\end{keypoints}

让我们先考虑一个单转子：一个质量为 $m$ 的粒子在圆周 $0\leqslant\theta<2\pi$ 上运动。其拉格朗日量为
\begin{equation*}
L=\frac{1}{2}m\dot\theta^{2}+K\cos\theta
\end{equation*}
哈密顿量具有形式
\begin{equation*}
H=\frac{(S^{z})^{2}}{2m}-K\cos\theta=\frac{(S^{z})^{2}}{2m}-\frac{K}{2}(a+a^{\dagger})
\end{equation*}
其中 $S^{z}$ 是角动量，而 $a=\mathrm{e}^{-\mathrm{i}\theta}$。如果选取基底态 $|n\rangle=(2\pi)^{-1/2}\mathrm{e}^{\mathrm{i}n\theta}$（$n$ 取整数），则我们发现 $S^{z}|n\rangle=n|n\rangle$，$a|n\rangle=|n-1\rangle$。

为了得到具有演生低能规范理论的最简单模型，让我们考虑以下由 $\theta_{\langle12\rangle}$、$\theta_{\langle23\rangle}$、$\theta_{\langle34\rangle}$ 和 $\theta_{\langle41\rangle}$ 描述的四转子模型（见图 10.10(a)）：
\begin{equation*}
H=\sum_{i}\left(U\big(S_{\langle i-1,i\rangle}^{z}-S_{\langle i,i+1\rangle}^{z}\big)^{2}+J\big(S_{\langle i,i+1\rangle}^{z}\big)^{2}\right), \tag{10.6.1}
\end{equation*}
其中我们已约定 $4+1\sim1$，$1-1\sim4$。希尔伯特空间由 $|n_{\langle12\rangle}n_{\langle23\rangle}n_{\langle34\rangle}n_{\langle41\rangle}\rangle$ 张成，其中整数 $n_{i,i+1}$ 是 $S_{\langle i,i+1\rangle}^{z}$ 的本征值。若 $U\gg J$，则低能激发由能量 $E=4Jn^{2}$ 的 $|nnnn\rangle$ 态描述。所有其余激发的能量均为 $U$ 的量级。

为了看清它与格点规范理论的联系，我们想写出四转子模型的拉格朗日量。如果把哈密顿量写成 $H=\frac{1}{2}P^{\top}VP$ 的形式，其中 $P^{\top}=(S_{\langle12\rangle}^{z},S_{\langle23\rangle}^{z},S_{\langle34\rangle}^{z},S_{\langle41\rangle}^{z})$，而
\begin{equation*}
V=\begin{pmatrix}
4U+2J & -2U & 0 & -2U\\
-2U & 4U+2J & -2U & 0\\
0 & -2U & 4U+2J & -2U\\
-2U & 0 & -2U & 4U+2J
\end{pmatrix},
\end{equation*}
则拉格朗日量为
\begin{equation*}
L=\frac{1}{2}\dot\Theta^{\top}M\dot\Theta,
\end{equation*}
其中 $\Theta^{\top}=(\theta_{\langle12\rangle},\theta_{\langle23\rangle},\theta_{\langle34\rangle},\theta_{\langle41\rangle})$，$M=V^{-1}$。

显然，在上述拉格朗日量中我们看不到规范理论的任何迹象。为了得到规范理论，我们需要换一种方法导出拉格朗日量。利用 $H$ 的路径积分表示，并注意到 $(\theta,S^{z})$ 是一对正则坐标--动量，我们发现
\begin{equation*}
Z=\int\mathcal{D}S^{z}\,\mathcal{D}\theta\;\mathrm{e}^{\mathrm{i}\int\mathrm{d}t\,\big(\sum_{i}S_{\langle i,i+1\rangle}^{z}\dot\theta_{\langle i,i-1\rangle}-H\big)}
\end{equation*}
引入场 $a_0(i)$（$i=1,2,3,4$）来解耦 $U$ 项，可以把上述路径积分改写为
\begin{equation*}
Z=\int\mathcal{D}S^{z}\,\mathcal{D}\theta\,\mathcal{D}a_0\;\mathrm{e}^{\mathrm{i}\int\mathrm{d}t\,\big(\sum_{i}S_{\langle i,i+1\rangle}^{z}\dot\theta_{\langle i,i+1\rangle}-\tilde H(S^{z},a_0)\big)}
\end{equation*}
其中 $\tilde H=\sum_{i}\Big(J(S_{\langle i,i+1\rangle}^{z})^{2}+a_0(i)\big(S_{\langle i-1,i\rangle}^{z}-S_{\langle i,i+1\rangle}^{z}\big)-\frac{a_0^{2}(i)}{4U}\Big)$。积掉 $S_{\langle i,i+1\rangle}^{z}$ 后，我们得到
\begin{equation*}
Z=\int\mathcal{D}\theta\,\mathcal{D}a_0\;\mathrm{e}^{\mathrm{i}\int\mathrm{d}t\,L(\theta,\dot\theta,a_0)}
\end{equation*}
其中拉格朗日量为
\begin{equation*}
L=\frac{1}{4J}\sum_{i}\left(\big(\dot\theta_{\langle i,i+1\rangle}+a_0(i)-a_0(i+1)\big)^{2}+\frac{a_0^{2}(i)}{4U}\right)
\end{equation*}
在大 $U$ 极限下，可以丢掉 $a_0^{2}(i)/4U$ 项，得到
\begin{equation*}
L=\frac{1}{4J}\sum_{i}\big(\dot a_{i,i+1}+a_0(i)-a_0(i+1)\big)^{2}
\end{equation*}
这正是单个方格上 $U(1)$ 格点规范理论的拉格朗日量，其中
\begin{equation*}
a_{i,i+1}=\theta_{\langle i,i+1\rangle}, \qquad a_{i+1,i}=-\theta_{\langle i,i+1\rangle}
\end{equation*}
作为格点规范场（见图 10.10(b)）。可以检验，上述拉格朗日量在下列变换下不变：
\begin{equation*}
a_{ij}(t)\to a_{ij}(t)+\phi_{j}(t)-\phi_{i}(t), \qquad a_0(i,t)\to a_0(i,t)+\dot\phi_{i}(t) \tag{10.6.2}
\end{equation*}
这个变换称为规范变换。我们注意到，低能波函数 $\Psi(a_{12},a_{23},a_{34},a_{41})$ 是 $|nnnn\rangle$ 态的叠加。所有低能态都是规范不变的，即在规范变换 $a_{ij}\to a_{ij}+\phi_{j}-\phi_{i}$ 下不变。

连续 $U(1)$ 规范理论的电场由 $\pmb{E}=\pmb{\dot a}-\partial a_{0}$ 给出（取 $c=e=1$ 单位）。在格点规范理论中，电场变成定义在链上的如下量：
\begin{equation*}
E_{ij}=\dot a_{ij}-\big(a_0(j)-a_0(i)\big)
\end{equation*}
我们看到，我们的格点规范拉格朗日量可以写成 $L=\frac{1}{4J}\sum_{i}E_{i,i+1}^{2}$。把它与连续 $U(1)$ 规范理论 $\mathcal{L}\propto\pmb{E}^{2}-\pmb{B}^{2}$ 相比较，可见我们的拉格朗日量只包含对应于 $\pmb{E}^{2}$ 的动能。更普遍的格点规范理论还包含对应于 $\pmb{B}^{2}$ 的势能项。

为了得到势能项，我们把转子模型推广为
\begin{equation*}
\begin{split}
H=\sum_{i}\Big(&U\big(S_{\langle i-1,i\rangle}^{z}-S_{\langle i,i+1\rangle}^{z}\big)^{2}+J\big(S_{\langle i,i+1\rangle}^{z}\big)^{2}\\
&+t\big(\mathrm{e}^{\mathrm{i}\theta_{\langle i-1,i\rangle}}\mathrm{e}^{\mathrm{i}\theta_{\langle i,i+1\rangle}}+\text{h.c.}\big)\Big),
\end{split} \tag{10.6.3}
\end{equation*}
我们注意到 $\langle nnn|\mathrm{e}^{\mathrm{i}\theta_{\langle i-1,i\rangle}}\mathrm{e}^{-\mathrm{i}\theta_{\langle i,i+1\rangle}}|nnn\rangle=0$。因此，在 $t$ 的一阶，新项对低能没有影响。新项的低能效应只在 $t$ 的二阶才出现。

如果带着新项重复上述计算，就得到以下拉格朗日量：
\begin{equation*}
L=\frac{1}{4J}\sum_{i}\left(\big(\dot a_{i,i+1}+a_0(i)-a_0(i+1)\big)^{2}-t\big(\mathrm{e}^{\mathrm{i}(a_{i-1,i}+a_{i,i+1})}+\text{h.c.}\big)+\frac{a_0^{2}(i)}{4U}\right)
\end{equation*}
在上述拉格朗日量中，为什么新项在 $t$ 一阶下没有低能效应，稍微难看出来一些。让我们集中考虑如下形式的涨落：\footnote{在格点规范理论中，这种类型的涨落被称为纯规范涨落。}
\begin{equation*}
a_{i-1,i}=\phi_{i}-\phi_{i-1}
\end{equation*}
积掉 $a_0(i)$ 之后，这类涨落的拉格朗日量具有形式 $L=\frac{1}{2}\phi_{i}m_{ij}\phi_{j}-\sum_{i}t\big(\mathrm{e}^{\mathrm{i}(\phi_{i-1}+\phi_{i+1})}+\text{h.c.}\big)$，其中 $m_{ij}=O(U^{-1})$。我们看到，在大 $U$ 极限下，上述形式的涨落是快速而强烈的。由于 $\phi_{i}$ 生活在紧致空间上（即 $\phi_{i}$ 与 $\phi_{i}+2\pi$ 表示同一个点），这些快速涨落全都具有 $U$ 量级的大能隙。现在我们看到，$t$ 项 $t\,\mathrm{e}^{\mathrm{i}(\phi_{i-1}-\phi_{i+1})}$ 由于 $\phi_{i}$ 的强涨落而平均为零，在 $t$ 一阶下没有任何影响。然而，在 $t$ 的二阶，会出现一项 $t^{2}\prod_{i=1,3}\mathrm{e}^{\mathrm{i}(\theta_{\langle i-1,i\rangle}+\theta_{\langle i,i+1\rangle})}=t^{2}\mathrm{e}^{\mathrm{i}\sum_{i}a_{i-1,i}}$。这样的项不依赖于 $\phi_{i}$，因而不会平均为零。于是我们预期低能有效拉格朗日量具有如下形式：
\begin{equation*}
L=\frac{1}{4J}\sum_{i}\left(\big(\dot a_{i,i+1}+a_0(i)-a_0(i+1)\big)^{2}+\frac{a_0^{2}(i)}{4U}\right)+g\cos\Phi \tag{10.6.4}
\end{equation*}
其中 $g=O(t^{2}/U)$，而 $\Phi=\sum_{i}a_{i,i+1}$ 是 $U(1)$ 规范场穿过方格的通量。上述拉格朗日量不依赖快速量子涨落，并且在规范变换 (10.6.2) 下不变。我们看到，正是 $a_{ij}$ 的某些组合的强量子涨落，导致了一个规范不变的低能有效拉格朗日量。

为了定量地计算 $g$，我们先要导出低能有效哈密顿量。如果把 $t$ 项当作微扰来处理，并把低能态当作简并态，那么在 $t$ 的二阶，我们有
\begin{align*}
&\langle n_{1},n_{1},n_{1},n_{1}|H_{eff}|nnnn\rangle\\
&=-\frac{t^{2}}{2U}\langle n_{1},n_{1},n_{1},n_{1}|\mathrm{e}^{\mathrm{i}(\theta_{\langle34\rangle}+\theta_{\langle41\rangle})}|n_{1},n_{1},n,n\rangle\langle n_{1},n_{1},n,n|\mathrm{e}^{\mathrm{i}(\theta_{\langle12\rangle}+\theta_{\langle23\rangle})}|nnnn\rangle\\
&\quad+\text{另外三个类似的项}\\
&=-\frac{2t^{2}}{U}
\end{align*}
其中 $n_{1}=n+1$。于是，低能有效哈密顿量为
\begin{equation*}
H=\sum_{i}\left(U\big(S_{\langle i-1,i\rangle}^{z}-S_{\langle i,i+1\rangle}^{z}\big)^{2}+J\big(S_{\langle i,i+1\rangle}^{z}\big)^{2}\right)-\frac{4t^{2}}{U}\cos(\Phi).
\end{equation*}
相应的拉格朗日量由式 (10.6.4) 给出，其中
\begin{equation*}
g=\frac{4t^{2}}{U}.
\end{equation*}
如前所述，纯规范涨落具有 $U$ 量级的大能隙。$U$ 以下的低能有效理论可以令 $U\to\infty$ 得到，
%FIG{10.11}: {(a) 三转子系统；(b) 相关的格点规范理论。}
并且我们得到
\begin{equation*}
L=\frac{1}{4J}\sum_{i}\big(\dot a_{i,i+1}+a_0(i)-a_0(i+1)\big)^{2}+g\cos\Phi \tag{10.6.5}
\end{equation*}

这个模型曾在 6.4 节中研究过。那里证明了，当 $g=0$ 时，能级为 $E_{n}=4Jn^{2}$，与式 (10.6.1) 的低能能级精确一致。因此，式 (10.6.1) 在低能下确实是一个规范理论。

这里我们想作一点评论。什么是规范理论？四转子模型在低能下真的是一个规范理论吗？规范理论的标准定义是含有规范势的理论。按照这个定义，如果我们坚持用规范势来写四转子模型的拉格朗日量，它就可以是一个规范理论。另一方面，四转子模型又不是规范理论，因为它的自然拉格朗日量不含有任何规范势。我们看到，规范理论并不是一个定义良好的物理概念。“一个模型是不是规范理论”并不是一个定义良好的问题。这几乎就像：如果我们把四个格点标记为 $(a,b,c,d)$，就管四转子模型叫规范理论；而如果我们把四个格点标记为 $(1,2,3,4)$，就管它叫非规范理论。从这个观点看，本节的讨论是相当形式化的，没有什么物理意义。

\subsection*{问题 10.6.1}

考虑式 (10.6.3) 所描述的 $i=1,2,3$ 三格点转子模型（见图 10.11）。

\begin{enumerate}
\item 导出低能格点规范理论的作用量。你只需把 $g$ 计算到 $O(1)$ 系数的精度。
\item 假定 $U\gg g\gg J$，计算低能下的能级。
\end{enumerate}

\subsection{量子转子格子与人工光}

\begin{keypoints}
\item 量子转子格子可以用一个格点 $U(1)$ 规范理论来描述。
\end{keypoints}

%FIG{10.12}: {转子格子、表示低能涨落的一个回路，以及一对电荷激发（A、B）。在大 $U$ 极限下，每个阴影方格四角处转子的角动量之和为零。}

\begin{keypoints}
\item 格点规范理论的解禁闭相包含人工光（artificial light）。
\end{keypoints}

在理解了只含少数几个转子的系统之后，我们就可以着手研究格子转子系统了。作为一个例子，我们将考虑一个正方格子，每条链上有一个转子（见图 10.12）。我们有
\begin{align*}
H=U\sum_{\pmb{i}}\Bigg(\sum_{\pmb{\alpha}}S_{\pmb{i}+\pmb{\alpha}}^{z}\Bigg)^{2}&+J\sum_{\pmb{i},\pmb{\alpha}}\big(S_{\pmb{i}+\pmb{\alpha}}^{z}\big)^{2}+\sum_{\pmb{i}}\Big(t_{1}\,\mathrm{e}^{\mathrm{i}\big(\theta_{\pmb{i}+\frac{\pmb{x}}{2}}-\theta_{\pmb{i}+\frac{\pmb{y}}{2}}\big)}\\
&+t_{2}\,\mathrm{e}^{\mathrm{i}\big(\theta_{\pmb{i}+\frac{\pmb{y}}{2}}-\theta_{\pmb{i}-\frac{\pmb{x}}{2}}\big)}+t_{3}\,\mathrm{e}^{\mathrm{i}\big(\theta_{\pmb{i}-\frac{\pmb{x}}{2}}-\theta_{\pmb{i}-\frac{\pmb{y}}{2}}\big)}+t_{4}\,\mathrm{e}^{\mathrm{i}\big(\theta_{\pmb{i}-\frac{\pmb{y}}{2}}-\theta_{\pmb{i}+\frac{\pmb{x}}{2}}\big)}+\text{h.c.}\Bigg) \tag{10.6.6}
\end{align*}
这里 $\pmb{i}=(i_x,i_y)$ 标记正方格子的格点，而 $\pmb{\alpha}=\pm\frac{\pmb{x}}{2},\pm\frac{\pmb{y}}{2}$。$U$ 项强制施加了“格点 $\pmb{i}$ 附近四个转子的总角动量为零”这一约束。该模型还可以推广到三维立方格子：
\begin{align*}
H=U\sum_{\pmb{i}}\Bigg(&\sum_{\pmb{\alpha}}S_{\pmb{i}+\pmb{\alpha}}^{z}\Bigg)^{2}+J\sum_{\pmb{i},\pmb{\alpha}}\big(S_{\pmb{i}+\pmb{\alpha}}^{z}\big)^{2}\\
&+\sum_{\pmb{i},\langle\pmb{\alpha}_1\pmb{\alpha}_2\rangle}\left(t_{\langle\pmb{\alpha}_1\pmb{\alpha}_2\rangle}\,\mathrm{e}^{\mathrm{i}(\theta_{\pmb{i}+\pmb{\alpha}_1}-\theta_{\pmb{i}+\pmb{\alpha}_2})}+\text{h.c.}\right) \tag{10.6.7}
\end{align*}
这里 $\pmb{i}=(i_x,i_y,i_z)$ 标记立方格子的格点，$\pmb{\alpha}=\pm\frac{\pmb{x}}{2},\pm\frac{\pmb{y}}{2},\pm\frac{\pmb{z}}{2}$，而两个标记 $\pmb{i}+\pmb{\alpha}_1$ 与 $\pmb{i}+\pmb{\alpha}_2$ 表示绕 $\pmb{i}$ 格点的两个最近邻转子。
